\documentclass[12pt]{article}
\usepackage{amsmath,amssymb,geometry,setspace,url}
\geometry{margin=1in}
\setstretch{1.2}
\begin{document}
\title{SKANN-SSL --- Ambient Noise Synthesis (Stage --1)}
\author{Oravont Systems LLP}
\date{}
\maketitle

SKANN Document D
Ambient Noise Synthesis (Stage )
Generating Time-Domain Waveforms from Spectral Models

Oravont Systems LLP

\section*{Introduction}
Document D describes the full methodology for generating time-domain ambient-noise waveforms  from the blended and digitised PSD  produced in Document B. This spectral-to-time conversion forms the basis of SKANN Stage , where synthetic ocean-noise waveforms are produced for:

pretraining the SSL backbone,

data augmentation,

controlled SNR mixing with vessel signatures,

benchmarking detection sensitivity.

Document D covers all required steps:

FFT grid construction,

amplitude spectrum computation,

random phase generation,

Hermitian symmetry enforcement,

inverse FFT,

amplitude normalisation,

multi-segment and multi-channel synthesis.

\section*{Overall Synthesis Pipeline}
The synthesis process maps a PSD to a time-domain signal using the following sequence:

Digitised PSD (from Document B):


Assign FFT frequency grid corresponding to sampling rate  and FFT size :


Interpolate  onto .

Compute amplitude spectrum:


Assign random phase:


Construct complex spectrum:


Impose Hermitian symmetry to ensure  is real.

Inverse FFT:


RMS scaling to match physical sea-state levels.

\section*{Diagram Placeholder: Complete Synthesis Flow}
Conceptual pipeline from PSD to time-domain waveform.

Construction of the FFT Frequency Grid

Let the chosen sampling rate and FFT size be:


The frequency resolution is:


Thus the -th discrete frequency is:


For real-valued signals, only


contain unique frequencies.

Interpolation of the PSD onto the FFT Grid

Document B produces a log-spaced set of PSD points:


where  may extend to 32 kHz.

To synthesise time-domain noise, the PSD must be evaluated on the FFT grid :


Interpolation is performed in log-frequency:


because ambient noise follows power-law relationships.

Amplitude Spectrum from PSD

A PSD expresses power per unit bandwidth:


\section*{Thus:}

This ensures that:


which matches the correct total mean-square pressure.

Worked Example: Sea State 3 at 1 kHz (Δf = 1 Hz)

Consider Sea State 3 at . From Document B, the noise level is:


Units: dB re 1 .

Convert this to linear PSD:


\section*{Numerically:}

\section*{With :}

Thus the one-sided FFT magnitude at 1 kHz is:


Time-Domain Interpretation of a Single Component

A positive-frequency bin and its conjugate form a real cosine with **twice** the one-sided amplitude:


\section*{Numerically:}

\section*{Random Phase Assignment}
Ambient noise has no deterministic phase, so each bin is assigned


and


Worked Example: Random Phase at 1 kHz

Amplitude from previous section:


Step 1: draw :


Step 2: convert to phase:


Step 3: construct the complex spectrum value:


\section*{Thus:}

\section*{Hermitian Symmetry: Explicit Conjugate Assignment}
For a real time-domain waveform, the FFT must satisfy:


With  and :


Compute the conjugate:


Assign to symmetric bin:


\section*{Resulting Time-Domain Tone}
The pair  produces:


\section*{At :}

\section*{At :}

\section*{Hermitian Symmetry Enforcement}
To guarantee real-valued samples:


and special bins:


After symmetry is imposed, the spectrum is ready for the inverse FFT.

Inverse FFT

The synthetic time-domain waveform is obtained via:


Because Hermitian symmetry is satisfied,  is guaranteed to be real:


The waveform retains the exact PSD shape defined in Document B and is ready for RMS-level scaling.

RMS Scaling and SPL Matching

The raw IFFT output has correct spectral shape but arbitrary amplitude. To match a desired sea-state or absolute level, the waveform is scaled to a target RMS value.

Let the generated signal have RMS:


Let the desired RMS correspond to a physical SPL:


with reference pressure .

Scale the signal:


This is applied once only, after the IFFT.

\section*{Generating Multi-Second Noise Segments}
The previous steps generate a single noise frame of length  samples. To produce long-duration noise (e.g., 2–30 seconds), multiple IFFTs are concatenated.

A frame produces duration:


\section*{Example:}

A multi-second duration  requires:


frames.

\section*{Avoiding Discontinuities Between Frames}
Simply concatenating IFFT outputs can produce discontinuities at frame boundaries because each frame uses independent random phase. SKANN uses one of two solutions:

\section*{Method 1: Overlap–Add (Recommended)}
Each frame is windowed with a smooth taper (50\% overlap):


Frames are placed into the output with:


The final signal is:


Window choice ensures:


so the energy is preserved.

\section*{Method 2: Correlated Phase Propagation}
To ensure smooth phase continuation, the random phases can be generated as:


where  is small and random. This preserves continuity at the cost of reduced randomness.

SKANN uses Method 1 because it is simpler, robust, and produces natural noise.

Long-Duration Synthesis (Minutes or Hours)

For extended-duration acoustic backgrounds (e.g., simulation of hydrophone streams):

use block-based OLA synthesis,

randomise seeds per block,

enforce continuity using Method 1,

periodically re-normalise RMS every 10–20 minutes.

Output lengths of 1 hour or more are straightforward.

\section*{Multi-Channel Synthesis (Optional)}
For hydrophone arrays, including vector sensors, SKANN supports multi-channel noise synthesis.

\section*{Independent-Channel Noise}
If channels are assumed uncorrelated:


with independent phases .

\section*{Spatially Correlated Noise}
If physical spatial correlations are required (e.g., HAVS dipole sensors):


where:

is the base noise spectrum,

is the channel’s spatial response.

\section*{Examples:}

These models allow SKANN to simulate:

tri-axial vector sensors,

hydrophone arrays,

correlated ambient-noise fields.

Diagram Placeholder: Multi-Segment OLA

Overlap–Add synthesis of long ambient-noise waveforms.

\section*{Complete Synthesis Algorithm (Pseudocode)}
This section summarises the full procedure in algorithmic form. The notation matches Documents B and C.

\section*{Input}
Digitised PSD  over 0–32 kHz (Document B)

Sampling rate  Hz

FFT size  (typically 32768 for  Hz)

Desired duration

Target RMS level

\section*{Output}
Time-domain ambient-noise waveform .

\section*{Pseudocode}
1. Construct FFT frequencies:
       f_k = k * f_s / N,       k = 0 .. N/2

2. Interpolate PSD onto FFT grid:
       Spp_fk = interp( Spp_fi → f_k )

3. Compute amplitude spectrum:
       A_k = sqrt( Spp_fk * Δf )

4. Generate random phases:
       φ_k ~ Uniform(0, 2π)

5. Construct complex spectrum:
       P[k] = A_k * exp(j φ_k)

6. Impose Hermitian symmetry:
       P[N-k] = conj(P[k])

7. Inverse FFT:
       p_frame = IFFT(P)

8. RMS-scale the frame:
       p_frame ← p_frame * (p_rms_target / p_rms_generated)

9. For long durations:
       - Generate M frames
       - Apply Hann window
       - Overlap-add with hop = N/2

10. Output:
       p_long[n]

Integration with Stage 0 (Document C)

When  is produced:

Resampling is not needed because synthesis already uses  Hz.

DC offset is already zero (expected from IFFT), but Stage 0 verifies:


Amplitude normalisation in Stage 0 is typically disabled for synthetic noise, since Stage  has applied the correct RMS.

Segmentation into fixed-length frames (e.g. 4096 samples) is performed exactly as for real hydrophone data.

Silence trimming is skipped; noise has no silent frames.

Thus the output of Document D becomes a valid input for Stage 0 and then for the learned filterbank in Stage 1.

Diagram Placeholder: Full Stage  to Stage 0 Interface

Integration of synthetic noise into the SKANN processing chain.

Appendix A — Unified List of Symbols (Documents A–D)

\section*{Spectral Quantities}
: FFT bin frequency [Hz]

: frequency resolution []

: pressure PSD [Pa/Hz]

: amplitude spectrum

: random phase

\section*{Time-Domain Quantities}
: discrete-time waveform

: RMS pressure

Pa: reference pressure

\section*{Noise Model Components}
: noise level [dB re 1 Pa/Hz]

SS: sea state (0–6)

: PSD scaling

\section*{Pandoc Compatibility Notes}
No TikZ or pgfplots used.

Figures implemented using figure with boxed placeholders.

Math expressions closed properly for DOCX conversion.

Standard packages only: amsmath, amssymb, geometry, setspace, url.

\end{document}